\sectionkicker{Source-backed technical notes}
\subsection*{スケールの意味が変わった――Chinchilla、PaLM、OPT/BLOOM、FlashAttention: Technical Notes}
本欄は記事本文で圧縮した一次資料上の情報を、比較・再検証しやすい形で整理したものである。時系列、確認済みの主張、留意点、一次資料URLを掲載する。
\smallskip
\noindent{\small\textit{Technical Notesでは、一次資料から確認した公開・提供・研究上の事実を記載する。数値・能力評価や提供元・著者による比較表現は、独立再現済みの結果ではなく帰属claimとして扱う。}}\par
\medskip
\subsection*{Theme at a glance}
\addcontentsline{toc}{subsection}{Theme at a glance}
\begin{center}
\footnotesize
\begin{tabularx}{\linewidth}{@{}p{0.20\linewidth}p{0.10\linewidth}p{0.14\linewidth}X@{}}
\toprule
資料 & 位置づけ & 種別 & 時系列 \\
\midrule
Chinchilla & 主要資料 & モデル & 2022-03-29, 2022-04-12 \\
PaLM & 主要資料 & モデル & 2022-04-04, 2022-04-05 \\
OPT-175B & 主要資料 & オープンウェイト & 2022-05-02, 2022-05-03 \\
BLOOM & 主要資料 & オープンウェイト & 2022-07-12, 2022-11-09 \\
FlashAttention & 主要資料 & 論文 & 2022-05-27 \\
GLM-130B & 補足資料 & オープンウェイト & 2022-10-05 \\
\bottomrule
\end{tabularx}
\end{center}
\normalsize

\begin{technicalnote}{Chinchilla}{主要資料}
\begingroup
% reader-facing Technical Notes break policy
\widowpenalty=10000
\clubpenalty=10000
\displaywidowpenalty=10000
\begin{tabularx}{\linewidth}{@{}>{\bfseries}p{0.22\linewidth}X@{}}
組織 & Google DeepMind \\
種別 & モデル \\
時系列 & 2022-03-29, 2022-04-12 \\
\end{tabularx}
\smallskip
{\bfseries 一次資料から整理したtechnical points}
\begin{itemize}[leftmargin=1.5em,itemsep=0.35em]
\item \textbf{一次情報で確認できる事実}: Google DeepMindの選定済み一次資料では、所与のtraining compute budgetの下でTransformer language modelのmodel sizeとtraining token数の配分を調べ、compute-optimal trainingでは両者を合わせて増やす必要があるというscaling relationを実験的に検討した。
\end{itemize}
\begin{minipage}{\linewidth}
% reader-facing Technical Notes generic boundary/source tail group
\begin{samepage}
{\bfseries 一次資料}
\begingroup\sloppy
\Urlmuskip=0mu plus 2mu\relax
\begin{itemize}[leftmargin=1.5em,itemsep=0.25em]
\item {\scriptsize\url{http://arxiv.org/abs/2203.15556v1}}
\item {\scriptsize\url{https://deepmind.google/blog/an-empirical-analysis-of-compute-optimal-large-language-model-training/}}
\end{itemize}
\endgroup
\end{samepage}
\end{minipage}
\endgroup
\end{technicalnote}
\medskip

\begin{technicalnote}{PaLM}{主要資料}
\begingroup
% reader-facing Technical Notes break policy
\widowpenalty=10000
\clubpenalty=10000
\displaywidowpenalty=10000
\begin{tabularx}{\linewidth}{@{}>{\bfseries}p{0.22\linewidth}X@{}}
組織 & Google Research \\
種別 & モデル \\
時系列 & 2022-04-04, 2022-04-05 \\
\end{tabularx}
\smallskip
{\bfseries 一次資料から整理したtechnical points}
\begin{itemize}[leftmargin=1.5em,itemsep=0.35em]
\item \textbf{一次情報で確認できる事実}: Google Researchの選定済み一次資料では、PaLMをdensely activated Transformer language modelとしてPathwaysにより複数TPU Podへまたがって学習し、model scaleがfew-shot learningへ与える影響を多数のlanguage understanding・generation taskで検討した。
\end{itemize}
\begin{minipage}{\linewidth}
% reader-facing Technical Notes generic boundary/source tail group
\begin{samepage}
{\bfseries 一次資料}
\begingroup\sloppy
\Urlmuskip=0mu plus 2mu\relax
\begin{itemize}[leftmargin=1.5em,itemsep=0.25em]
\item {\scriptsize\url{http://arxiv.org/abs/2204.02311v5}}
\item {\scriptsize\url{https://research.google/blog/pathways-language-model-palm-scaling-to-540-billion-parameters-for-breakthrough-performance/}}
\end{itemize}
\endgroup
\end{samepage}
\end{minipage}
\endgroup
\end{technicalnote}
\medskip

\begin{technicalnote}{OPT-175B}{主要資料}
\begingroup
% reader-facing Technical Notes break policy
\widowpenalty=10000
\clubpenalty=10000
\displaywidowpenalty=10000
\begin{tabularx}{\linewidth}{@{}>{\bfseries}p{0.22\linewidth}X@{}}
組織 & Meta AI \\
種別 & オープンウェイト \\
時系列 & 2022-05-02, 2022-05-03 \\
\end{tabularx}
\smallskip
{\bfseries 一次資料から整理したtechnical points}
\begin{itemize}[leftmargin=1.5em,itemsep=0.35em]
\item \textbf{一次情報で確認できる事実}: Meta AIの選定済み一次資料では、OPTは125Mから175B parameterまでのdecoder-only pretrained Transformer群として提示され、研究者がlarge language modelのweightsを直接調査できるaccessを主眼の一つに置いた。 論文と同時期の公開では、model weightsだけでなく、学習時のinfrastructure上の課題を記録したlogbookと実験用codeも研究再現性のための構成要素として扱われた。
\end{itemize}
\begin{minipage}{\linewidth}
% reader-facing Technical Notes generic boundary/source tail group
\begin{samepage}
{\bfseries 一次資料}
\begingroup\sloppy
\Urlmuskip=0mu plus 2mu\relax
\begin{itemize}[leftmargin=1.5em,itemsep=0.25em]
\item {\scriptsize\url{http://arxiv.org/abs/2205.01068v4}}
\item {\scriptsize\url{https://ai.meta.com/blog/democratizing-access-to-large-scale-language-models-with-opt-175b/}}
\end{itemize}
\endgroup
\end{samepage}
\end{minipage}
\endgroup
\end{technicalnote}
\medskip

\begin{technicalnote}{BLOOM}{主要資料}
\begingroup
% reader-facing Technical Notes break policy
\widowpenalty=10000
\clubpenalty=10000
\displaywidowpenalty=10000
\begin{tabularx}{\linewidth}{@{}>{\bfseries}p{0.22\linewidth}X@{}}
組織 & BigScience / Hugging Face \\
種別 & オープンウェイト \\
時系列 & 2022-07-12, 2022-11-09 \\
\end{tabularx}
\smallskip
{\bfseries 一次資料から整理したtechnical points}
\begin{itemize}[leftmargin=1.5em,itemsep=0.35em]
\item \textbf{一次情報で確認できる事実}: BigScience / Hugging Faceの選定済み一次資料では、BLOOMはROOTS corpusで学習したdecoder-only Transformer language modelとして記述され、46の自然言語と13のprogramming languageを含むmultilingual corpusを用い、modelとcodeをResponsible AI Licenseの下で公開した。
\end{itemize}
\begin{minipage}{\linewidth}
% reader-facing Technical Notes generic boundary/source tail group
\begin{samepage}
{\bfseries 一次資料}
\begingroup\sloppy
\Urlmuskip=0mu plus 2mu\relax
\begin{itemize}[leftmargin=1.5em,itemsep=0.25em]
\item {\scriptsize\url{http://arxiv.org/abs/2211.05100v4}}
\item {\scriptsize\url{https://huggingface.co/blog/bloom}}
\end{itemize}
\endgroup
\end{samepage}
\end{minipage}
\endgroup
\end{technicalnote}
\medskip

\begin{technicalnote}{FlashAttention}{主要資料}
\begingroup
% reader-facing Technical Notes break policy
\widowpenalty=10000
\clubpenalty=10000
\displaywidowpenalty=10000
\begin{tabularx}{\linewidth}{@{}>{\bfseries}p{0.22\linewidth}X@{}}
組織 & authors \\
種別 & 論文 \\
時系列 & 2022-05-27 \\
\end{tabularx}
\smallskip
{\bfseries 一次資料から整理したtechnical points}
\begin{itemize}[leftmargin=1.5em,itemsep=0.35em]
\item \textbf{一次情報で確認できる事実}: authorsの選定済み一次資料では、FlashAttentionはGPU memory hierarchy上のread/write量を明示的に考慮するIO-awareなexact attention algorithmとして提案された。 algorithmはtilingによってHBMとon-chip SRAMの間のmemory accessを削減し、近似attentionへ置き換えずにself-attentionのsystems bottleneckへ対処する。
\end{itemize}
\begin{minipage}{\linewidth}
% reader-facing Technical Notes generic boundary/source tail group
\begin{samepage}
{\bfseries 一次資料}
\begingroup\sloppy
\Urlmuskip=0mu plus 2mu\relax
\begin{itemize}[leftmargin=1.5em,itemsep=0.25em]
\item {\scriptsize\url{http://arxiv.org/abs/2205.14135v2}}
\end{itemize}
\endgroup
\end{samepage}
\end{minipage}
\endgroup
\end{technicalnote}
\medskip

\begin{technicalnote}{GLM-130B}{補足資料}
\begingroup
% reader-facing Technical Notes break policy
\widowpenalty=10000
\clubpenalty=10000
\displaywidowpenalty=10000
\begin{tabularx}{\linewidth}{@{}>{\bfseries}p{0.22\linewidth}X@{}}
組織 & Tsinghua / authors \\
種別 & オープンウェイト \\
時系列 & 2022-10-05 \\
\end{tabularx}
\smallskip
{\bfseries 一次資料から整理したtechnical points}
\begin{itemize}[leftmargin=1.5em,itemsep=0.35em]
\item \textbf{一次情報で確認できる事実}: Tsinghua / authorsの選定済み一次資料では、GLM-130Bを英語・中国語のbilingual pre-trained language modelとして提示し、loss spikeやdivergenceへの対処を含むtraining design・efficiency・stability上のengineeringを記録するとともに、model weights、code、training logs、toolkitを公開した。
\end{itemize}
\begin{minipage}{\linewidth}
% reader-facing Technical Notes generic boundary/source tail group
\begin{samepage}
{\bfseries 一次資料}
\begingroup\sloppy
\Urlmuskip=0mu plus 2mu\relax
\begin{itemize}[leftmargin=1.5em,itemsep=0.25em]
\item {\scriptsize\url{http://arxiv.org/abs/2210.02414v2}}
\end{itemize}
\endgroup
\end{samepage}
\end{minipage}
\endgroup
\end{technicalnote}
\medskip
